%% Essuie-glace : dessin d'ensemble (pièces numérotées, pièces coupées hachurées).
\documentclass[tikz,border=4pt]{standalone}
\usepackage{liaisons}
%% Dessin d'ensemble simplifié de l'essuie-glace de voiture (système à quatre barres).
%% Inclus par mecanisme-essuie-glace-dessin.tex (pièces numérotées, en noir) et
%% mecanisme-essuie-glace-classes.tex (pièces colorées par classe d'équivalence) ;
%% chacun définit \piece{<n° de pièce>}{<chemin>} (l'apparence d'une pièce) et \eglabels.
%% Repère : x horizontal, y vertical, z vers l'observateur. Cotes en cm.
%% Points caractéristiques : O (0,0) axe moteur, A (0.8,0) maneton de la manivelle,
%% D (4,0) axe du balancier, B (4,2.2) articulation bielle/balancier.

%% \barre{P}{Q}{demi-largeur} : chemin fermé du rectangle d'axe PQ (pas de rotation de repère,
%% les quatre sommets sont construits avec la bibliothèque calc).
\newcommand\barre[3]{($(#1)!#3!90:(#2)$) -- ($(#2)!#3!-90:(#1)$)
                  -- ($(#2)!#3!90:(#1)$) -- ($(#1)!#3!-90:(#2)$) -- cycle}

\newcommand\essuieglace{%
  \begin{scope}[line width=0.6pt, line join=round]
    \coordinate (O) at (0,0);    \coordinate (A) at (0.8,0);
    \coordinate (D) at (4,0);    \coordinate (B) at (4,2.2);
    \coordinate (T) at (4,3.05);
    % ---- 0 : platine, boulonnée sur la carrosserie (semelle + montant du palier D)
    \piece{0}{(-1.55,-2.05) rectangle (5.45,-1.55)}
    \piece{0}{(3.42,-1.55) -- (3.42,-0.35) arc (180:0:0.58) -- (4.58,-1.55) -- cycle}
    % ---- 6 : carter du motoréducteur + palier de l'arbre de sortie en O
    \piece{6}{(-1.3,-1.55) rectangle (1.0,-0.42)}
    \piece{6}{(O) circle (0.42)}
    % ---- 3 : balancier (levier calé sur l'axe D)
    \piece{3}{\barre{D}{B}{0.17cm}}
    \piece{3}{(D) circle (0.3)}
    \piece{3}{(B) circle (0.2)}
    % ---- 4 : bras d'essuie-glace (prolonge le balancier au-delà de B)
    \piece{4}{\barre{B}{T}{0.11cm}}
    % ---- 5 : balai (monture + lame de caoutchouc), perpendiculaire au bras
    \piece{5}{(2.95,2.93) rectangle (5.05,3.19)}
    \piece{5}{(3.05,2.72) rectangle (4.95,2.93)}
    % ---- 2 : bielle
    \piece{2}{\barre{A}{B}{0.16cm}}
    \piece{2}{(A) circle (0.2)}
    \piece{2}{(B) circle (0.2)}
    % ---- 1 : manivelle (plateau de sortie du motoréducteur) et son maneton en A
    \piece{1}{\barre{O}{A}{0.19cm}}
    \piece{1}{(O) circle (0.28)}
    \piece{1}{(A) circle (0.14)}
    % ---- points caractéristiques
    \foreach \p/\n/\lab in {O/{(0,0)}/{(-0.62,0.30)}, A/{(0.8,0)}/{(1.10,0.48)},
                            B/{(4,2.2)}/{(3.56,2.46)}, D/{(4,0)}/{(4.80,0.20)}} {
      \draw[blue, line width=0.5pt] \n ++(-0.13,0) -- ++(0.26,0) \n ++(0,-0.13) -- ++(0,0.26);
      \node[blue, font=\small, inner sep=1pt] at \lab {\p};
    }
    \eglabels
  \end{scope}}

\begin{document}
\begin{tikzpicture}[x=1cm,y=1cm]
% une apparence par pièce : 0 et 6 sont coupées (hachures), les autres sont vues de face
\newcommand\piece[2]{\ifcase#1\relax
  \hachures{#2}%                                     0 : platine
  \or \draw[line width=0.6pt, line join=round, fill=white] #2;%   1 : manivelle
  \or \draw[line width=0.6pt, line join=round, fill=white] #2;%   2 : bielle
  \or \draw[line width=0.6pt, line join=round, fill=white] #2;%   3 : balancier
  \or \draw[line width=0.6pt, line join=round, fill=white] #2;%   4 : bras
  \or \draw[line width=0.6pt, line join=round, fill=white] #2;%   5 : balai
  \or \hachuresb{#2}%                                6 : carter du motoréducteur
  \fi}
\newcommand\eglabels{%
  \foreach \n/\pt/\lab in {0/{(1.9,-1.55)}/{(1.9,-2.75)}, 6/{(-1.3,-1.0)}/{(-2.3,-1.0)},
                           1/{(0.4,0.19)}/{(0.15,1.15)}, 2/{(2.4,1.1)}/{(1.55,1.75)},
                           3/{(4.17,1.3)}/{(5.0,1.0)}, 4/{(4.11,2.6)}/{(5.0,2.35)},
                           5/{(5.05,3.06)}/{(5.9,3.5)}} {
    \draw[line width=0.4pt] \pt -- \lab;
    \node[fill=white, inner sep=1.5pt, font=\small\bfseries] at \lab {\n};
  }}
\essuieglace
\repere{(-2.4,-2.6)}{x}{y}{1}
\end{tikzpicture}
\end{document}
